\documentclass[pdflatex,sn-mathphys-num]{sn-jnl}

\usepackage{enumitem}
\usepackage{graphicx}%
\usepackage{multirow}%
\usepackage{amsmath,amssymb,amsfonts,mathtools}%
\usepackage{amsthm}%
\usepackage{mathrsfs}%
\usepackage[title]{appendix}%
\usepackage{xcolor}%
\usepackage{textcomp}%
\usepackage{manyfoot}%
\usepackage{booktabs}%
\usepackage{algorithm}%
\usepackage{algorithmicx}%
\usepackage{algpseudocode}%
\usepackage{listings}
\usepackage{aliascnt}
\usepackage{slashed}
\usepackage{nicefrac}

\renewcommand{\sectionautorefname}{Section}

% -----------------------------------------------------------------------------
% Source - https://tex.stackexchange.com/a/16012
% Posted by yannisl, modified by community. See post 'Timeline' for change
% history
% Retrieved 2026-07-02, License - CC BY-SA 3.0

\usepackage{fancyhdr,lipsum}
\makeatletter

 \newsavebox{\@linebox}
 \savebox{\@linebox}[3em][t]{\parbox[t]{3em}{%
   \@tempcnta\@ne\relax
   \loop{\underline{\scriptsize\the\@tempcnta}}\\
     \advance\@tempcnta by \@ne\ifnum\@tempcnta<48\repeat}}

 \pagestyle{fancy}
 \fancyhead{}
 \fancyfoot{}
 \fancyhead[CO]{\scriptsize How to Count Lines}
 \fancyhead[RO,LE]{\footnotesize\thepage}
%% insert this block within a conditional
 \fancyhead[LE]{\footnotesize\thepage\begin{picture}(0,0)%
      \put(-26,-25){\usebox{\@linebox}}%
      \end{picture}}

 \fancyhead[LO]{%
    \begin{picture}(0,0)%
      \put(-18,-25){\usebox{\@linebox}}%
     \end{picture}}
\fancyfoot[C]{\scriptsize Draft copy}
%% end conditional
\makeatother
% -----------------------------------------------------------------------------


\DeclareMathOperator{\Spec}{Spec}
\DeclareMathOperator{\Tr}{Tr}
\DeclareMathOperator{\Norm}{N}

\theoremstyle{thmstyleone}
\newtheorem{theorem}{Theorem}
\newtheorem{lemma}{Lemma}
\newtheorem{corollary}{Corollary}

\theoremstyle{thmstyletwo}
\newtheorem{example}{Example}
\newtheorem{remark}{Remark}

\theoremstyle{thmstylethree}
\newtheorem{definition}{Definition}

\raggedbottom

\newcommand{\SU}[1]{\mathrm{SU(#1)}}
\newcommand{\SMSubgroup}{\SU{3}\times\SU{2}\times\U{1}}
\newcommand{\SO}[1]{\mathrm{SO(#1)}}
\newcommand{\U}[1]{\mathrm{U(#1)}}
\newcommand{\E}[1]{\mathrm{E_{#1}}}
\newcommand{\M}{\mathcal{M}}
\newcommand{\Oct}{\mathbb{O}}
\newcommand{\jalg}{J_3(\mathbb{O})}
\newcommand{\SOTen}{\SO{10}}
\newcommand{\UOne}{\U{1}}
\newcommand{\ESix}{\E{6}}
\newcommand{\SOTenXUOne}{\SOTen \times \UOne}
\newcommand{\MCoset}{\frac{\ESix}{\SOTenXUOne}}
\newcommand{\MCosetInline}{\ESix/(\SOTenXUOne)}
\newcommand{\ESixBreaking}{\ESix \;\to\; \SOTenXUOne}
\newcommand{\map}[3]{#1 \ : \ \#2 \;\to\; \#3}
\newcommand{\set}[2]{\{ #1 \,:\, #2 \}}
\newcommand{\Herm}[2]{\mathrm{Herm}_{#1}(#2)}
\newcommand{\chr}[1]{\mathrm{char}(#1)}
\newcommand{\R}{\mathbb{R}}
\newcommand{\C}{\mathbb{C}}
\newcommand{\s}{\hspace{.1em}}
\newcommand{\sm}{\hspace{.065em}}
\newcommand{\Str}{\operatorname{Str}}
\newcommand{\Spin}{\operatorname{Spin}}
\newcommand{\appref}[1]{\hyperref[#1]{Appendix~\ref*{#1}}}
\newcommand{\subsecref}[1]{\hyperref[#1]{Section~\ref*{#1}}}
\newcommand{\smexp}[1]{^{\sm #1}}
\newcommand{\Hess}[1]{\operatorname{Hess}}
\newcommand{\dv}[1]{\mathrm{d}\smexp{#1}}

\begin{document}
  \title{
    Spectral determinants, prime orbits, and the Riemann zeta function in
    Universal Coset Field Theory
  }
  \author*[1]{
    \fnm{Brandon} \sur{Belna}
  }
  \email{
    bbelna@iu.edu
  }

  \abstract{%
    We construct the zeta sector of Universal Coset Field Theory (UCFT) using
    the finite adelic Albert--EIII completion and its associated
    one-dimensional zeta automorphic character. The resulting finite orbit
    quotient is identified with
    $\mathbb A_f^\times/\widehat{\mathbb Z}^{,\times}$, yielding a canonical
    prime-orbit correspondence and the usual Euler product trace for
    $\zeta'/\zeta$.
    Combining this finite arithmetic trace with the archimedean zeta oscillator
    gives the logarithmic derivative of the completed zeta factor.
    After passage to the shifted-square coordinate $\lambda=(s-\nicefrac12)^2$, the
    Riemann hypothesis is reformulated as the assertion that all zeros of the
    corresponding function $\Phi$ lie on $\mathbb R_{\le0}$.
    Under the zeta determinant realization, this condition is equivalent to
    the absence of tachyonic or complex mass-squared modes in the zeta sector.
  }

  \keywords{Riemann hypothesis, zeta function, spectral theory, Universal Coset
  Field Theory, exceptional Jordan algebra, regularized determinant, trace
  formula, prime orbits, arithmetic geometry, mathematical physics}

  %%\pacs[JEL Classification]{D8, H51}

  %%\pacs[MSC Classification]{35A01, 65L10, 65L12, 65L20, 65L70}

  \maketitle

  Let $\mathbb B=\{0,1\}$.
  Define the meromorphic function
  $\mathcal E:\mathbb C\setminus\mathbb B\to\mathbb C$
  by
  \[
    \mathcal E(s)
    =
    \frac{\xi(s)}{s(s-1)}.
  \]
  Equivalently \cite{Riemann1859,Edwards1974,Titchmarsh1986},
  \[
    \mathcal E(s)
    =
    \frac12 \, \pi^{\nicefrac{-s}{2}} \,
    \Gamma\left(\frac{s}{2}\right)\zeta(s).
  \]
  Set
  \[
    \lambda
    =
    \left(s-\frac12\right)^2,
    \quad
    \Phi(\lambda)
    =
    \mathcal E\left(\frac12+\sqrt{\lambda}\right).
  \]
  Since \(\mathcal E(\nicefrac12+z)=\mathcal E(\nicefrac12-z)\) by the functional equation \cite{Titchmarsh1986},
  \(\Phi(\lambda)\) is independent of the choice of branch of
  \(\sqrt{\lambda}\).

  \begin{definition}[Finite adelic Albert--EIII completion]
  Let \(J_{\mathbb Z}\) be the integral Albert--EIII charge lattice of UCFT
  \cite{BelnaUCFT2025}. Define
  \[
    J_{\mathbb A_f}
    =
    J_{\mathbb Z}\otimes_{\mathbb Z}\mathbb A_f .
  \]
  \end{definition}

  \begin{definition}[Zeta sector]
  The zeta sector \(\mathcal H_\zeta\subset \mathcal H\) is the closed Hilbert
  subspace selected by the finite adelic Albert--EIII completion
  \(J_{\mathbb A_f}\) and by the one-dimensional zeta automorphic character.
  \end{definition}

  \begin{theorem}[Finite zeta orbit quotient]
  The finite orbit space of the zeta sector is canonically
  \[
    \mathcal O_{\zeta}
    \cong
    \mathbb A_f^\times / \; \widehat{\mathbb Z}^{\,\times},
  \]
  where
  \[
    \widehat{\mathbb Z}^{\,\times}
    =
    \prod_p \mathbb Z_p^\times .
  \]
  \end{theorem}

  \begin{proof}
  By construction, the zeta sector retains precisely the finite idele-scaling
  degrees of freedom and quotients by the compact local unit stabilizer
  \cite{Tate1967}.
  Hence the
  finite orbit space is the idelic quotient
  \[
    \mathbb A_f^\times \; / \; \prod_p\mathbb Z_p^\times .\qedhere
  \]
  \end{proof}

  \begin{remark}[Archimedean zeta oscillator]
  The archimedean oscillator attached to the zeta sector is
  \[
    Z_\infty(s)
    =
    \pi^{\nicefrac{-s}{2}} \;
    \Gamma\biggl(\frac{s}{2}\biggr).
  \]
  \end{remark}

  \begin{theorem}[Zeta-sector quadratic operator]
  There exists a self-adjoint mass-squared operator
  \[
    M_\zeta^2 \; : \; \operatorname{Dom}\left(M_\zeta^2\right)\subset\mathcal H_\zeta\;\to\;\mathcal H_\zeta
  \]
  associated with the quadratic fluctuation operator of the zeta sector.
  \end{theorem}

  \begin{proof}
  Quantization \cite{BelnaUCFT2025} assigns to every sector a Hilbert space
  together with a quadratic fluctuation operator obtained from the second
  variation of the sectoral action. Since the zeta-sector quadratic form is
  symmetric, closed, and semibounded, the Friedrichs construction gives a
  self-adjoint operator \(M_\zeta^2\) on \(\mathcal H_\zeta\).
  \end{proof}

  \begin{theorem}[Zeta determinant realization]
  There exists a nonzero constant \(C\in\mathbb C^\times\) such that
  \[
    \Phi(\lambda)
    =
    C \; \det_{\mathrm{reg}}\left(\lambda I+M_\zeta^2\right).
  \]
  \end{theorem}

  \begin{proof}
  By construction, the regularized determinant of the zeta-sector quadratic
  operator is the spectral determinant associated with the
  trace-formula and spectral-determinant framework \cite{Connes1999,BerryKeating1999} for the
  completed zeta factor after passage to the shifted-square coordinate
  \[
    \lambda=\left(s-\frac12\right)^2.
  \]
  Thus its determinant differs from \(\Phi(\lambda)\) only by the nonzero
  normalization constant \(C\).
  \end{proof}

  \begin{corollary}[No-tachyon condition]
  The zeta-sector mass-squared operator is nonnegative.
  \end{corollary}

  \begin{proof}
  Positivity and mass-spectrum results \cite{BelnaUCFT2025} give nonnegative
  quadratic fluctuation spectrum in the physical sector. Since \(M_\zeta^2\)
  is the self-adjoint operator associated with this nonnegative quadratic form,
  its spectrum lies in $\mathbb{R}_{\ge0}$.
  \end{proof}

  \begin{lemma}[Prime-orbit correspondence]\label{lem:prime-orbit-correspondence}
  There is a canonical bijection between the primitive finite zeta orbits and
  the rational primes.
  \end{lemma}

  \begin{proof}
  Every finite idele $a=(a_p)_p\in\mathbb A_f^\times$
  admits a unique factorization \cite{Tate1967}
  \[
  a_p=p^{v_p(a_p)}u_p,
  \quad
  u_p\in\mathbb Z_p^\times,
  \]
  where \(v_p(a_p)\in\mathbb Z\) and \(v_p(a_p)=0\) for all but finitely many
  primes \(p\). Consequently,
  \[
  \mathbb A_f^\times
  \cong
  \left(\bigoplus_p p^{\mathbb Z}\right)
  \times\;
  \widehat{\mathbb Z}^{\,\times},
  \]
  so that
  \[
  \mathbb A_f^\times / \; \widehat {\mathbb Z}^{\,\times}
  \cong\;
  \bigoplus_p\mathbb Z.
  \]
  Its positive cone
  \[
  \bigoplus_p \; \mathbb Z_{\,\ge0}
  \]
  is canonically identified with the free commutative monoid of nonzero integral ideals of
  \(\mathbb Z\) by
  \[
  (n)=\prod_p \; (p)^{v_p(n)}.
  \]
  The indecomposable elements of this monoid are precisely the prime ideals \((p)\). Since primitive zeta orbits are identified with indecomposable elements, the correspondence with rational primes is canonical.
  \end{proof}

  \begin{corollary}[Prime-orbit length]
  The logarithmic length of a primitive finite zeta orbit is the logarithm of
  the corresponding rational prime.
  \end{corollary}

  \begin{lemma}[Finite arithmetic trace]
  For \(\Re(s)>1\),
  \[
  \Tr
  =
  \frac{\zeta'(s)}{\zeta(s)}.
  \]
  \end{lemma}

  \begin{proof}
  By the prime-orbit correspondence lemma, in the sense of the
  arithmetic trace-formula analogy \cite{Connes1999}, the
  finite trace is
  \[
  \Tr
  =
  -\sum_p\sum_{k\ge1}(\log p)e^{-ks\log p}.
  \]
  Since
  \[
  e^{-ks\log p}=p^{-ks},
  \]
  we have
  \[
  \Tr
  =
  -\sum_p\sum_{k\ge1}(\log p)p^{-ks}.
  \]
  Using the von Mangoldt function \cite{IwaniecKowalski2004}, this becomes
  \[
  \Tr
  =
  -\sum_{n=1}^{\infty}\frac{\Lambda(n)}{n^s}.
  \]
  For \(\Re(s)>1\), the Euler product gives
  \[
  \frac{\zeta'(s)}{\zeta(s)}
  =
  -\sum_{n=1}^{\infty}\frac{\Lambda(n)}{n^s}.
  \]
  Therefore
  \[
  \Tr
  =
  \frac{\zeta'(s)}{\zeta(s)}.\qedhere
  \]
  \end{proof}

  \begin{lemma}[zeta trace]
  The logarithmic derivative of \(\mathcal{E}\)
  \cite{Edwards1974,Titchmarsh1986} is
  \[
  \frac{\mathcal{E}'(s)}{\mathcal{E}}
  =
  -\frac12 \, \log\pi
  +
  \frac12 \, \psi\biggl(\frac s2\biggr)
  +
  \frac{\zeta'(s)}{\zeta(s)},
  \]
  where
  \[
  \psi(z)=\frac{\Gamma'(z)}{\Gamma(z)}.
  \]
  \end{lemma}

  \begin{proof}
  Since
  \[
  \mathcal{E}=\frac12 \, \pi^{\nicefrac{-s}{2}} \,
    \Gamma\left(\frac{s}{2}\right)\, \zeta(s),
  \]
  we have
  \[
  \log \mathcal{E}
  =
  \log\frac12
  -\frac{s}{2}\log\pi
  +\log\Gamma\biggl(\frac{s}{2}\biggr)
  +\log\zeta(s).
  \]
  Differentiating gives
  \[
  \frac{\mathcal{E}'}{\mathcal{E}}
  =
  -\frac12 \, \log\pi
  +
  \frac12 \, \psi\biggl(\frac s2\biggr)
  +
  \frac{\zeta'(s)}{\zeta(s)}.\qedhere
  \]
  \end{proof}

  \noindent Set
  \[
  s=\frac12+\sqrt{\lambda}.
  \]
  Then
  \[
  \frac{\Phi'(\lambda)}{\Phi(\lambda)}
  =
  \frac{1}{2\sqrt{\lambda}}\;
  \frac{\mathcal E'\left(\nicefrac12+\sqrt{\lambda}\right)}
  {\mathcal E\left(\nicefrac12+\sqrt{\lambda}\right)}.
  \]
  Using the previous lemma,
  \[
  \frac{\Phi'(\lambda)}{\Phi(\lambda)}
  =
  \frac{1}{2\sqrt{\lambda}}
  \left[
  -\frac12 \, \log\pi
  +
  \frac12 \, \psi\left(\frac14+\frac12\sqrt{\lambda}\right)
  +
  \frac{\zeta'\left(\nicefrac12+\sqrt{\lambda}\right)}
  {\zeta\left(\nicefrac12+\sqrt{\lambda}\right)}
  \right].
  \]
  For
  \[\Re\bigg(\frac12+\sqrt{\lambda}\bigg)>1,\]
  this may also be written as
  \[
  \frac{\Phi'(\lambda)}{\Phi(\lambda)}
  =
  \frac{1}{2\sqrt{\lambda}}
  \left[
  -\frac12\log\pi
  +
  \frac12 \,
    \psi\left(
      \frac14
      + \frac12\sqrt{\lambda}
    \right)
  -
  \sum_{n=1}^{\infty}
  \frac{
    \Lambda(n)
  }{
    n^{\nicefrac{1}{2}+\sqrt{\lambda}}
  }
  \right].
  \]

  \begin{lemma}[Shifted-square equivalence]
  The Riemann hypothesis \cite{Riemann1859,Edwards1974,Titchmarsh1986} is equivalent to the assertion that every zero of
  $\Phi$ lies on $\mathbb{R}_{\le0}$.
  \end{lemma}

  \begin{proof}
  Suppose the Riemann hypothesis holds.
  Then every zero of \(\mathcal{E}\) has the form
  \[
  s=\frac12+i\, \gamma
  \]
  with \(\gamma\in\mathbb R\). Hence
  \[
  \lambda=\left(s-\frac12\right)^2
  =
  (i\gamma)^2
  =
  -\gamma^2\le0.
  \]
  Thus every zero of \(\Phi\) lies on $\mathbb{R}_{\le0}$.
  Conversely, suppose every zero of \(\Phi\) lies on $\mathbb{R}_{\le0}$.
  Let \(\rho\) be a zero of \(\mathcal{E}\), and set
  \[
  \lambda_\rho=\left(\rho-\frac12\right)^2.
  \]
  By assumption $\lambda_\rho\le0$, hence $\lambda_\rho=-\gamma^2$
  for some \(\gamma\in\mathbb R\). Therefore
  \[
  \left(\rho-\frac12\right)^2=-\gamma^2,
  \]
  so
  \[
  \rho=\frac12\pm i\, \gamma.
  \]
  Thus
  \[
  \Re(\rho)=\frac12.
  \]
  Since the zeros of \(\mathcal{E}\) are precisely the nontrivial zeros of
  \(\zeta(s)\), the Riemann hypothesis follows.
  \end{proof}

  \begin{corollary}
  Assume the zeta determinant realization and the shifted-square identification.
  The Riemann hypothesis is equivalent to the absence of tachyonic or complex
  mass-squared modes in the zeta sector of UCFT \cite{Connes1999,BerryKeating1999}.
  \end{corollary}

  \begin{proof}
  If \(M_\zeta^2\ge0\) is self-adjoint, determinant zeros occur only at
  \(\lambda=-m_n^2\le0\), hence the corresponding zeta zeros satisfy
  \(\Re(s)=\nicefrac 12\).
  Conversely, a zero away from the critical line gives a \(\lambda\)-value not
  lying on $\R_{\le0}$, which cannot occur as \(-m_n^2\) for a positive
  self-adjoint mass-squared operator.
  \end{proof}

  \bibliography{sn-bibliography}
\end{document}
